\documentclass[10pt]{article}
\usepackage[margin=2.2cm]{geometry}
\usepackage{amsmath,amssymb,amsthm,bm}
\usepackage[hidelinks]{hyperref}
\newtheorem{proposition}{Proposition}
\newtheorem{corollary}{Corollary}
\theoremstyle{definition}
\newtheorem{assumption}{Assumption}
\newcommand{\E}{\mathbb{E}}
\newcommand{\tr}{\operatorname{tr}}
\title{Spatial Discriminability in Reverberant Rooms\\[2pt]\large Theory note}
\author{D. Bassett}
\date{Working draft, \today}

\begin{document}
\maketitle

\noindent This note sets out the observation model, the cue statistics, and the closed-form results on which the paper rests. The anechoic metric of~\cite{bassett2026} is the limit of no reverberant energy. Code: \texttt{sdr/}; verification: \texttt{tests/test\_model.py}.

\section{Observation model}

A static source at direction $\theta\in S^2$ and distance $r_s$ radiates a stationary Gaussian signal of power spectral density $P(f)$ to a listener whose head is described by directional transfer functions $\bm h(f,\phi)=[h_1,h_2]^{\top}$. In a room the binaural transfer function is a sum over arrivals $k$ (direct path $k=0$, image sources $k\ge1$) with amplitudes $a_k$, delays $\tau_k$ and directions $\phi_k$, all of which depend on $\theta$:
\begin{equation}
\bm H(f,\theta)=\sum_k a_k(f)\,e^{-j2\pi f\tau_k}\,\bm h(f,\phi_k).
\end{equation}
The auditory system analyses the ear signals in bands $b$ with power responses $|G_b(f)|^2$. The expected band cross-spectral matrix is
\begin{equation}
\bm R_b(\theta)=\int |G_b(f)|^2 P(f)\,\bm H\bm H^{\mathsf H}\,df
=\sum_{k,l} a_k a_l^{*}\int |G_b|^2 P\,\bm h_k\bm h_l^{\mathsf H}e^{-j2\pi f(\tau_k-\tau_l)}\,df .
\label{eq:Rb}
\end{equation}

\begin{proposition}[Coherence window]\label{prop:window}
If $P$, $a_k$ and $\bm h_k$ vary slowly across band $b$, the pair $(k,l)$ contributes $a_ka_l^{*}\,\bm h_k\bm h_l^{\mathsf H}\,e^{-j2\pi f_b\Delta\tau}\,\kappa_b(\Delta\tau)\,S_b$ to $\bm R_b$, where $\Delta\tau=\tau_k-\tau_l$, $S_b=\int|G_b|^2P\,df$, and for a fourth-order gammatone with bandwidth $\beta_b=1.019\,\mathrm{ERB}(f_b)$
\begin{equation}
\kappa_b(\Delta\tau)=e^{-s}\Big(1+s+\tfrac{2}{5}s^2+\tfrac{1}{15}s^3\Big),\qquad s=2\pi\beta_b|\Delta\tau|.
\end{equation}
$\kappa_b$ falls to one half at $|\Delta\tau|\approx0.45/\mathrm{ERB}(f_b)$: about 6~ms at 500~Hz, 3.4~ms at 1~kHz, and 1~ms at 4~kHz.
\end{proposition}
\begin{proof}
$|G_b(f)|^2=(1+((f-f_b)/\beta_b)^2)^{-4}$; its Fourier transform is the Mat\'ern-$7/2$ covariance above, modulated by $e^{-j2\pi f_b\Delta\tau}$.
\end{proof}

Arrivals separated by less than the coherence window add coherently within a band and shape the cue \emph{means}; arrivals separated by more add in power and reduce interaural \emph{coherence}. The boundary between early and late energy is therefore a property of the auditory filter, $0.45/\mathrm{ERB}(f)$, and not a fixed time such as 50 or 80~ms. When the late field is dense and isotropic, its incoherent sum tends to the diffuse-field matrix of the same head,
\begin{equation}
\bm\Gamma(f)=\int_{S^2}\bm h(f,\phi)\bm h(f,\phi)^{\mathsf H}\,\frac{d\phi}{4\pi},
\end{equation}
so that, with source power normalised to one and an omnidirectional direct-to-reverberant ratio $\mathrm{DRR}=1/\eta$,
\begin{equation}
\bm R(f,\theta)=\bm h_e(f,\theta)\bm h_e(f,\theta)^{\mathsf H}+\eta\,\bm\Gamma(f),
\label{eq:R}
\end{equation}
where $\bm h_e$ is the direct path plus the coherent early arrivals. The pilot below uses $\bm h_e=\bm h(\cdot,\theta)$.

\section{Cues and their statistics}

The cues are those of the Barumerli observer~\cite{barumerli2023}, written as functionals of $\bm R$: band levels $L_{i,b}=5\log_{10}S_{i,b}$ with $S_{i,b}=\int|G_b|^2R_{ii}\,df$ and their positive one-ERB gradients; the broadband level difference $\mathrm{ILD}=10\log_{10}(\int R_{11}/\int R_{22})$; and the interaural time difference $\tau^{*}$ at the peak of the envelope of the low-passed cross-correlation, $|\int_0^\infty L(f)R_{12}(f)e^{j2\pi f\tau}df|$, on the Reijniers JND scale. The factor 5 in $L_{i,b}$ reproduces the half-wave-rectified level of~\cite{barumerli2023}, whose gradient energy it matches to within 11\% (Section~\ref{sec:check}).

\begin{assumption}[Internal noise as decorrelation]\label{ass:noise}
Internal noise on the interaural cues acts as an internal coherence ceiling $\rho\in(0,1)$: an interaural cue estimated from signals of coherence $\gamma$ behaves as if their coherence were $\rho\gamma$. Its variance is the anechoic variance $\sigma^2$ of~\cite{bassett2026} times the ratio of estimator variances at coherence $\rho\gamma$ and $\rho$.
\end{assumption}

Coherence is taken per auditory band, $\gamma_b=R_{12,b}/\sqrt{R_{11,b}R_{22,b}}$ with $R_{ij,b}=\int|G_b|^2R_{ij}\,df$ over a 1-ERB gammatone bank from 80~Hz, because a deterministic set of arrivals has unit coherence at every frequency and loses it only by averaging within a band (Proposition~\ref{prop:window}); Rakerd and Hartmann likewise specify coherence per 1/3-octave band. Each band contributes in proportion to its bandwidth $B_b$. The Knapp--Carter bound for the delay and the delta method for the level ratio give
\begin{align}
v_{\mathrm{ITD}}&=\sigma_{\mathrm{ITD}}^2\,\frac{I(\gamma^{\mathrm{an}})}{I(\gamma)},\qquad I(\gamma)=\sum_b L(f_b)\,B_b\,\omega_b^2\,\frac{\rho^2\gamma_b^2}{1-\rho^2\gamma_b^2},\\
v_{\mathrm{ILD}}&=\sigma_{\mathrm{ILD}}^2\,\frac{V(\bm R)}{V(\bm R_{\mathrm{an}})},\qquad
V(\bm R)=\sum_b\frac{1}{B_b}\Big[\tfrac{R_{11,b}^2}{S_1^2}+\tfrac{R_{22,b}^2}{S_2^2}-\tfrac{2\rho^2|R_{12,b}|^2}{S_1S_2}\Big],
\end{align}
with $S_i=\sum_bR_{ii,b}$ and the anechoic reference taken at the same direction.
The monaural cues keep their anechoic variance $\sigma_{\mathrm{mon}}^2$: for a Gaussian signal the relative fluctuation of a band power depends on the bandwidth--duration product, not on the room. The sample-size factor is common to numerator and denominator, so the duration of the stimulus does not enter; $\rho$ is the only free parameter.

\subsection{Fisher metric}
For independent Gaussian cues with means $\mu_m(\theta)$ and variances $v_m(\theta)$,
\begin{equation}
G_{ij}(\theta)=\sum_m\frac{\partial_i\mu_m\,\partial_j\mu_m}{v_m}+\frac12\sum_m\partial_i\log v_m\,\partial_j\log v_m .
\label{eq:G}
\end{equation}
\begin{proposition}[Anechoic limit]
As $\eta\to0$, $\gamma\to1$, $v_m\to\sigma_m^2$, the second term vanishes, and $G$ is the anechoic metric.
\end{proposition}

\section{Closed-form consequences of the diffuse model}

Write $D_{i,b}(\theta)=S^{\mathrm{dir}}_{i,b}(\theta)/(\eta\,\Gamma_{ii,b})$ for the \emph{local} DRR of ear $i$ in band $b$: direct over diffuse power after head filtering.

\begin{proposition}[Level compression]\label{prop:mon}
$\partial L^{\mathrm{room}}_{i,b}=\dfrac{D_{i,b}}{1+D_{i,b}}\,\partial L^{\mathrm{an}}_{i,b}$, and each spectral gradient combines the compression factors of its two bands.
\end{proposition}
\begin{proof}
$\eta\Gamma_{ii,b}$ does not depend on $\theta$, so $\partial\log(S^{\mathrm{dir}}+\eta\Gamma)=\partial S^{\mathrm{dir}}/(S^{\mathrm{dir}}+\eta\Gamma)$.
\end{proof}
\begin{corollary}[Notch filling]
Bands where the direct spectrum is low, at pinna notches and at the shadowed ear, have the lowest local DRR and lose spectral-cue information first. Elevation acuity is set by the local DRR in the notch region, not by the broadband DRR.
\end{corollary}

\begin{proposition}[ILD compression]\label{prop:ild}
$\partial\,\mathrm{ILD}^{\mathrm{room}}=\frac{10}{\ln10}\big(\alpha_1\partial\log S_1^{\mathrm{dir}}-\alpha_2\partial\log S_2^{\mathrm{dir}}\big)$ with broadband $\alpha_i=D_i/(1+D_i)$, and $|\mathrm{ILD}^{\mathrm{room}}|\le10\log_{10}(1+D_{\mathrm{near}})+\mathrm{ILD}_{\mathrm{diffuse}}$.
\end{proposition}
For a lateral source the far ear is shadowed, its local DRR is small, and its contribution to the ILD slope vanishes first: ILD is compressed most for lateral sources, as measured by Shinn-Cunningham et al.~\cite{shinn2005}.

\begin{proposition}[Interaural coherence]\label{prop:coh}
For a source in the median plane of a left--right symmetric head, $|h_1|=|h_2|$ and $\Gamma_{11}=\Gamma_{22}$, so
\begin{equation}
\gamma=\frac{D\,e^{j\varphi}+\gamma_d}{D+1},
\end{equation}
with $\gamma_d=\Gamma_{12}/\Gamma_{11}$ the diffuse-field coherence of the head and $\varphi$ the interaural phase of the direct path.
\end{proposition}
At low frequencies $\gamma_d\to1$ and reverberation leaves the coherence high but pulls the interaural phase towards zero; at high frequencies $\gamma_d\to0$ and coherence falls as $D/(D+1)$. The ITD cue therefore loses information through two mechanisms with opposite frequency dependence: bias at low frequencies and variance at high.

\begin{proposition}[Critical DRR]
The ITD variance doubles relative to anechoic when $\rho^2|\gamma|^2/(1-\rho^2|\gamma|^2)=\tfrac12\rho^2/(1-\rho^2)$ over the ITD band, i.e.\ at $|\gamma|^{2}=1/(2-\rho^2)$. With Proposition~\ref{prop:coh} and $\gamma_d\approx0$ this gives $D^{*}=|\gamma^*|/(1-|\gamma^*|)$: about 18~dB for the fitted $\rho=0.984$ (Section~\ref{sec:c1}).
\end{proposition}

\section{Predictions}
\begin{enumerate}
\item Thresholds in azimuth and elevation rise smoothly with falling DRR, with a knee near the critical DRR; elevation thresholds depend on the local DRR of the notch bands and not on interaural coherence.
\item Azimuth acuity is lost faster for lateral than for frontal sources at moderate DRR, through ILD compression at the shadowed ear.
\item The weight of the cues shifts with DRR: frontally, the ITD share of the information falls and the spectral share rises.
\item Reflections within $0.45/\mathrm{ERB}(f)$ of the direct sound bias cues rather than decorrelate them; the window shortens with frequency.
\item Acuity maps across a room follow the local DRR and interaural coherence at each position, not reverberation time alone.
\end{enumerate}
The envelope ITD of the Barumerli cue set is biased outward by the diffuse field: the diffuse interaural cross-correlation is a plateau over the range of ITDs of the head, and the envelope of a band-limited plateau peaks at its edges. At $-5$~dB DRR the envelope peak of a $90^\circ$ source moves from 696 to 730~\textmu s while the fine-structure peak moves from 704 to 702~\textmu s, in line with the robustness of the ITD within $\pm1$~ms reported in~\cite{shinn2005}. The bias is a property of the estimator, not a prediction.

\section{Calibration (C1): the internal coherence ceiling}\label{sec:c1}
Rakerd and Hartmann~\cite{rakerd2010} measured lateralisation of 1/3-octave noise bands recorded in rooms, with interaural coherence 0.2--1.0 and level differences removed. Their psychometric functions were read from the vector figures of the article (Figs.~2, 3 and 8; 72 points, axis calibration residual $<0.01$ percentage points; curve assignment checked against the article's curve labels by overlay, \texttt{figures/c1\_overlay.png}; \texttt{data/rakerd\_hartmann\_2010.csv}). A single-interval judgement of the sign of an ITD $\Delta t$ is modelled as
\begin{equation}
P(\text{consistent})=\Phi\!\left(\frac{\Delta t}{\sigma_0\sqrt{m(\gamma)}}\right),\qquad
m(\gamma)=\frac{(1-\rho^2\gamma^2)/(\rho^2\gamma^2)}{(1-\rho^2)/\rho^2},
\end{equation}
with a scale $\sigma_0$ per band and the waveform coherence $\gamma$ for the 225 and 715~Hz bands, the envelope coherence $\gamma^2$ for the 2850~Hz band (\texttt{experiments/c1\_calibrate.py}). One shared ceiling fits all three bands: $\rho=0.984$ (approximately $0.978$--$0.989$, treating each point as 300 binomial judgements), with root-mean-square errors of 1.3, 4.2 and 4.8 percentage points. Separate ceilings per band do not improve the likelihood. The 715 and 2850~Hz bands, tested only up to coherence 0.8, do not constrain $\rho$ on their own; the 225~Hz band, tested at 0.85 and 1.0, does. The 715~Hz data of a second listener group (Fig.~8b), not used in the fit, are predicted to 4.0 percentage points.

The greater sensitivity of the 225~Hz band to decorrelation, which Rakerd and Hartmann emphasise, is carried by the scale and not by the ceiling: $\sigma_0=48$~\textmu s at 225~Hz against 10~\textmu s at 715~Hz, a ratio of 4.6 where the $\omega^2$ weighting of the Knapp--Carter information predicts 3.2. The model overpredicts performance at coherence 0.2 in the 715~Hz band (74\% at 175~\textmu s against 60\% measured), where listeners are close to chance.

\section{Pilot (E1): direct path plus diffuse field}
Forty-one SONICOM heads, $\rho=0.984$ with its interval 0.978--0.989, DRR from $\infty$ to $-10$~dB (\texttt{experiments/e1\_drr.py}). The median threshold at $d'=1$ in azimuth rises from $1.6^\circ$ to $2.3^\circ$, $5.2^\circ$ and $22^\circ$ at 10, 0 and $-10$~dB for a frontal source ($5.0^\circ$--$5.4^\circ$ at 0~dB over the interval of $\rho$), and from $7.6^\circ$ to $9.8^\circ$, $15.0^\circ$ and $42^\circ$ at $90^\circ$. In elevation at the front it rises from $4.0^\circ$ to $4.8^\circ$, $8.2^\circ$ and $31^\circ$, independently of $\rho$. Frontally the share of the information carried by the ITD falls from 32\% to 7\% at 0~dB, while the spectral share rises from 14\% to 35\%.

Recomputing the metric with room means and anechoic variances (compression only), and with anechoic means and room variances (inflation only), separates the two mechanisms (\texttt{experiments/e4\_mechanism.py}, 20 heads). For a frontal source $\sqrt{\det G}$ falls to 0.59 of its anechoic value at 10~dB: 0.82 from compression and 0.74 from inflation, whose product, 0.60, matches the full model. Inflation saturates near 0~dB once the interaural cues have lost their information (0.37--0.75 of anechoic, depending on direction), whereas compression keeps falling (0.40 at 0~dB, 0.09 at $-10$~dB). At $90^\circ$ inflation has no effect and the loss is compression alone. Above about 5~dB the room acts mainly by making the interaural cues noisier; below 0~dB, mainly by biasing all cues.

\section{Precedence (E2): temporal weighting}
\paragraph{The stationary model has no precedence.} For a source and a single copy delayed by more than about 1~ms, the stationary model carries the same information about the source's direction whether it leads or lags: for a frontal target with a copy at $45^\circ$, and for $30^\circ$ with $-30^\circ$, the ratio of the information about the target as lead to that as lag is 0.99--1.03 at 2 and 10~ms. Below about 1~ms the summed cross-correlation envelope is bimodal and the envelope ITD jumps between peaks between neighbouring directions, so the Gaussian cue model does not apply there.

\paragraph{Weighted variant.} Litovsky and Shinn-Cunningham~\cite{litovsky2001} measured ITD jnds of the lead and of the lag of two 1-ms noise bursts under headphones, at delays of 1--15~ms, with the ITD of the other burst roved over $\pm500$~\textmu s (six listeners). Their Fig.~3 was digitised from the scanned article (78 points, checked by overlay; \texttt{data/litovsky\_shinn\_2001\_fig3.csv}). The target burst's ITD is estimated against the other burst, which dilutes its effective interaural coherence within the binaural integration window,
\begin{equation}
\gamma=\frac{a_t}{a_t+a_n\,c(D)},\qquad
w(D)=1-e^{-D/T_s},\qquad c(D)=e^{-D/T_b},
\end{equation}
with $(a_t,a_n)=(1,w)$ for the lead and $(w,1)$ for the lag, and $\mathrm{jnd}=A_s\sqrt{m(\gamma;\rho)}$ with $\rho=0.984$ carried over from C1 and a scale $A_s$ per listener (the single-burst control lacks the roving ITD and does not set the scale). Log jnds have Gaussian error; the 26 jnds at or above the adaptive-track ceiling (320~\textmu s) are right-censored (\texttt{experiments/e2\_precedence.py}). The fit gives $T_s=5.6$~ms (profile interval 4.6--7.1~ms) and $T_b=50$~ms, with a residual standard deviation of 0.28 in log jnd, and is preferred over the stationary model ($w\equiv1$) by an AIC difference of 93. The weight of a later arrival recovers to 0.16 at 1~ms, 0.30 at 2~ms and 0.83 at 10~ms. Freeing $\rho$ improves the AIC by 4.6 but drives it to 0.63 and $T_b$ to infinity; the burst experiment constrains the coherence ceiling poorly, and $\rho$ is kept at its calibrated value. The model misses the late recovery of S2's lag discrimination and the continued rise of lead jnds beyond 10~ms for three listeners.

\paragraph{Consequence for rooms.} With $T_s\approx6$~ms, reflections arriving within a few milliseconds of the direct sound carry a fraction of their weight in the cues, while reflections later than about 15~ms carry nearly all of it. The diffuse late field of E1 is therefore essentially unaffected by the weighting; early reflections bias the cues less than the stationary model predicts.

\section{Shoebox room (E1b): image sources and listener position}
An image-source model of a shoebox room supplies the arrivals (\texttt{sdr/ism.py}): images to third order are summed coherently, each filtered by the DTF of its arrival direction, and the reverberant energy they do not carry is added as the diffuse field of the head, so that the total matches the statistical reverberant energy $16\pi r^2/R$ with room constant $R$. DTFs at arbitrary directions are interpolated barycentrically on the convex hull of the measurement grid after removing each ear's delay (phase slope 200--1500~Hz); a virtual vertex at the south pole carries the mean of the $-45^\circ$ ring, so floor reflections from below the measured region are approximated by that ring. Interpolation is exact at the measured directions and errs by 0.5~dB (median) and 2.3~\textmu s in leave-one-out tests.

In an $8\times6\times3$~m room with absorption 0.1 and 0.3 (Eyring reverberation times 1.2 and 0.4~s; DRR $-3.1$ and 2.8~dB for a source at 0.9~m), twelve heads, three listener positions (centre; 1~m from a side wall, right ear towards it; 1~m from a corner) (\texttt{experiments/e1b\_shoebox.py}):
\begin{itemize}
\item At the same DRR the image-source room degrades acuity less than the diffuse model: frontal azimuth thresholds of $7.1^\circ$ against $8.1^\circ$ at absorption 0.1 and $3.3^\circ$ against $4.1^\circ$ at 0.3, elevation $10.6^\circ$ against $11.5^\circ$ and $6.5^\circ$ against $7.0^\circ$. Early images fall within the coherence window of the low bands and keep the ear signals more coherent (frontal band coherence 0.55 against 0.38, and 0.88 against 0.67). The diffuse model is therefore conservative by 10--25\% in threshold.
\item Near the wall and in the corner, frontal coherence falls (0.45 against 0.55 at absorption 0.1, 0.73 against 0.88 at 0.3), as Shinn-Cunningham et al.~\cite{shinn2005} measured, and frontal thresholds rise modestly ($7.3^\circ$--$7.5^\circ$ against $7.1^\circ$, and $3.5^\circ$--$3.9^\circ$ against $3.2^\circ$). Near the boundaries the images carry more energy than the statistical reverberant energy, so the nominal DRR overstates the true one there.
\item For sources on the wall side the threshold at $90^\circ$ is lower near the wall ($15.6^\circ$--$17.2^\circ$ against $18.6^\circ$ at the centre, absorption 0.1). Such sources lie 0.1~m from the wall and their first-order image nearly coincides with them, acting as a stronger direct path.
\item Precedence weighting changes thresholds by up to $0.4^\circ$, except for wall-side sources, where it discounts the near-coincident image and raises the threshold at $90^\circ$ by $1.1^\circ$--$1.4^\circ$; the floor reflection, 4.8~ms after the direct sound, keeps a weight of 0.58.
\end{itemize}

\section{Measured rooms (E3): the late field against measured BRIRs}\label{sec:e3}
The Surrey real-room set~\cite{hummersone2010} holds BRIRs of a Cortex Mk.2 head and torso simulator in four rooms (A, office, $T_{60}=0.32$~s; B, classroom, 0.47~s; C, lecture theatre, 0.68~s; D, seminar room, 0.89~s), for a loudspeaker at 1.5~m and ear height, $-90^\circ$ to $90^\circ$ in $5^\circ$ steps, with pseudo-anechoic responses of the same head (\texttt{sdr/brir.py}, \texttt{experiments/e3\_measured.py}). Each response is cut where its energy envelope meets the noise floor and split by a window of $[-1, 2.5]$~ms about the onset, with a 1-ms fall, into a direct part and its complement; the direct part is the anechoic reference, and agrees with the pseudo-anechoic set to 2--6\% in threshold. Cross-spectra are formed at the full response length and summed into the bins of the HRTF grid, and every response of a room is divided by the common transfer function of its direct parts. The DRR measured this way is 5.9--9.6~dB (median over directions). With only the horizontal arc measured, information is computed along azimuth by local quadratic regression over $\pm20^\circ$; on the SONICOM horizontal plane this estimate agrees with the order-6 SH metric to a median of 4\% in threshold anechoically and 8--10\% at 10 to 0~dB DRR, and to 5--7\% in the room-to-anechoic ratio (twelve heads).

The measured responses are compared with two models that keep the direct part and replace reverberant energy by the room's own late field: the cross-spectral matrix after 80~ms, averaged over the 37 directions, scaled to the measured energy at each direction.
\begin{itemize}
\item \emph{Direct part plus diffuse field} at the measured DRR, the model of E1, over-predicts thresholds at 73\% of directions, by a median of 12\% (median absolute error 17\%). The error is largest at the front (12--41\% by room) and small or reversed laterally ($-9$ to $+25\%$). Measured frontal band coherence is higher than the model's (0.75--0.85 against 0.62--0.78): early reflections fall within the coherence window of the low bands, as in the image-source rooms of E1b, where the diffuse model was conservative by 10--25\%.
\item \emph{The first 80~ms as measured plus the diffuse field} reproduces the measured thresholds to a median of 3\% (90th percentile 9\%), with correlations across direction of 0.96--0.99 in log threshold. The late field is therefore well described as a direction-independent diffuse field at its measured energy: its band coherence matches that of a head in an isotropic field (0.56--0.88 at 200~Hz against 0.71 for the SONICOM median, below 0.3 above 500~Hz in both), although its level varies with source position by 7--10~dB across the arc. In these rooms only 2--18\% of the reverberant energy at 1.5~m arrives after 80~ms, so most of the loss of acuity is carried by early reflections.
\end{itemize}
Reverberation raises the measured azimuth threshold by a median factor of 1.8--2.6 over the direct part. Precedence weighting ($T_s=5.6$~ms) lowers thresholds within $\pm45^\circ$ by 7--12\% and leaves the comparison unchanged (median errors 13\% and 3\% for the two weighted models); near $\pm60^\circ$, where the measured threshold has a local maximum, it raises them by up to 16\%. The test concerns the room acoustics of the model, not its psychophysics: no discrimination thresholds have been measured in these rooms.

\section{Redundancy (E5): acuity against room-acoustic parameters}
Three shoebox rooms ($5\times4\times3$, $8\times6\times3$ and $14\times10\times5$~m), five absorptions (0.05--0.5), three random listener positions per room, sources at 1~m (ears at 1.2~m) and, in the tallest room, at 2~m (ears at 1.8~m), seven heads and eight source directions: 3\,360 source--receiver pairs in 420 configurations (\texttt{experiments/e5\_redundancy.py}). The metric is computed as elsewhere (order-6 SH fits, here over 176 near-uniform source directions above $-45^\circ$) for the room and for the anechoic case, and the acuity loss at a source direction is $\tfrac12\log(\det G_{\mathrm{an}}/\det G)$. DRR, C80, LF (5--80~ms, $\cos^2$ weighting), IACC (full response, mean of the 500, 1000 and 2000~Hz octaves) and $T_{60}$ are computed from the same arrivals. Ordinary least squares with folds that leave out whole configurations gives the cross-validated variance explained (\texttt{experiments/e5\_analysis.py}):
\begin{center}
\begin{tabular}{lccc}
Predictors & Overall & Azimuth & Elevation\\\hline
DRR & 0.80 & 0.71 & 0.81\\
DRR, C80, LF, IACC, $\log T_{60}$, linear & 0.81 & 0.72 & 0.82\\
the same, with squares and products & 0.85 & 0.77 & 0.85\\
local DRR (DRR plus head gain), with its square & 0.83 & 0.76 & 0.81\\
quadratic set plus head gain & 0.93 & 0.86 & 0.92\\
quadratic set plus source-direction indicators & 0.96 & 0.92 & 0.93\\
\end{tabular}
\end{center}
At the level of a seat, the loss of acuity in these rooms is largely redundant with DRR; C80, LF, IACC and $T_{60}$ add five points beyond it. What the parameters cannot carry is the dependence on source direction: 18\% of the variance lies across directions within a configuration, and at the same DRR the threshold area grows by a factor of 1.7--2.3 more for sources at $150^\circ$ and $180^\circ$ than at the front, and by 0.6 times as much at $90^\circ$. The head's direct-to-diffuse gain at each direction, which turns the omnidirectional DRR into the local DRR of the theory, accounts for most of that dependence: rear sources are attenuated by the pinna at high frequencies, and a lateral source is boosted at the near ear. The local DRR alone predicts acuity loss as well as all five standard parameters together.

Shoebox rooms with uniform absorption and a statistical late field make the standard parameters strongly correlated (correlation with the loss: DRR $-0.89$, IACC $-0.81$, C80 $-0.78$, LF $0.78$, $\log T_{60}$ $0.71$), which favours redundancy; rooms with non-uniform absorption, scattering and focusing, measured or simulated in detail, are needed before the conclusion can be generalised.

\section{Consistency with the anechoic metric}\label{sec:check}
With $\eta=0$ the spectral-domain cues reproduce the time-domain cues of~\cite{bassett2026} for nine SONICOM heads (\texttt{experiments/e0\_anechoic.py}): ILD and ITD correlate at $1.000$ with gradient energy within 2\%; the spectral gradients correlate at 0.86--0.90 with gradient energy within 11\%. $\sqrt{\det G}$ agrees within 5\% and tensor shape to a median affine-invariant distance of 0.12--0.15.

\begin{thebibliography}{9}
\bibitem{bassett2026} D. Bassett, ``Evaluating spatial discriminability preservation in HRTF upsampling using Fisher information,'' 2026.
\bibitem{barumerli2023} R. Barumerli, P. Majdak, M. Geronazzo, D. Meijer, F. Avanzini, and R. Baumgartner, ``A Bayesian model for human directional localization of broadband static sound sources,'' \emph{Acta Acustica}, vol. 7, 12, 2023.
\bibitem{knapp1976} C. H. Knapp and G. C. Carter, ``The generalized correlation method for estimation of time delay,'' \emph{IEEE Trans. Acoust., Speech, Signal Process.}, vol. 24, no. 4, pp. 320--327, 1976.
\bibitem{shinn2005} B. G. Shinn-Cunningham, N. Kopčo, and T. J. Martin, ``Localizing nearby sound sources in a classroom: Binaural room impulse responses,'' \emph{J. Acoust. Soc. Am.}, vol. 117, no. 5, pp. 3100--3115, 2005.
\bibitem{litovsky2001} R. Y. Litovsky and B. G. Shinn-Cunningham, ``Investigation of the relationship among three common measures of precedence: Fusion, localization dominance, and discrimination suppression,'' \emph{J. Acoust. Soc. Am.}, vol. 109, no. 1, pp. 346--358, 2001.
\bibitem{hummersone2010} C. Hummersone, R. Mason, and T. Brookes, ``Dynamic precedence effect modeling for source separation in reverberant environments,'' \emph{IEEE Trans. Audio, Speech, Lang. Process.}, vol. 18, no. 7, pp. 1867--1871, 2010. BRIRs: \texttt{github.com/IoSR-Surrey/RealRoomBRIRs}.
\bibitem{rakerd2010} B. Rakerd and W. M. Hartmann, ``Localization of sound in rooms. V. Binaural coherence and human sensitivity to interaural time differences in noise,'' \emph{J. Acoust. Soc. Am.}, vol. 128, no. 5, pp. 3052--3063, 2010.
\end{thebibliography}
\end{document}
